% \begin{savequote}
% \includegraphics[width=4cm]{./images/gabriel-christine-aass/hygiene-hello-002-edited.jpg}\tiny\sffamily
% 
% Keime mit unterschiedlichen
% Farben auf einem Nährboden.
% \\
% Bild: \emph{Marina Marinkovic}
% \end{savequote}
\def\ChapterCode{HYG}
\chapter
[\CO{[HYG]} Hygiene]
{Hygiene}
\label{chp:HYG}

\ChapterIntro
{%
~~
}
{%
\begin{description}
  \item[Maintainer:] Sebastian Gabriel
  \item[Autoren:] Diverse
  \item[Reviewer:] Alexander Blacky, Standard-Reviewprozess
  \item[Version:] {\DocVersion} ({\DocVersionInfo})
  \item[SHA1:] {\DocGitCommit}
\end{description}

{\TxTeilkapitelAASS}
}

\nocite{Longmore2007}
\nocite{IntroClinEmergMed4}
\nocite{Boecker2}
\nocite{Fauci2008}
\nocite{ForthPharma8}

% \Changelog{
% Changelog:
% 
% \begin{description}
%     \item[2009-02-20] Fertig für Kurs 2009/02.
%     \item[2009-02-25] Start Changelog.
%     \item[2009-04-02] Übernahme aus Openoffice in \LaTeX
% \end{description}
% }

\clearpage % TEMP temp clearpage
\section{Grundlagen}

\subsection{Wichtige Begriffe}

%\begin{multicols}{2}
\begin{labeling}[:]{Inkubationszeitx}
  \item[\T{Infektion}]\index{Infektion}
  Eindringen von Krankheitserregern in den Organismus.
  \item[\T{Pathogenität}]\index{Pathogenität} 
  Fähigkeit eines Krankheitserregers oder Stoffs eine
  Krankheit oder einen Schaden auszulösen.
  \item[\T{Resistenz}]\index{Resistenz}
  \begin{inparadesc}
    \item[Allg.:] Widerstandsfähigkeit.
    \item[Spez.:]
    \begin{inparaenum}
      \item W. eines Organismus gegen die
      krankmachenden Eigenschaften eines Erregers.
      \item W. eines Krankheitserregers gegen eine Behandlung (z.\,B.
      durch Antibiotika, Desinfektionsmittel, \ldots{}).
    \end{inparaenum}
  \end{inparadesc}
  \item[\T{Inkubationszeit}]\index{Inkubationszeit}
  symptomlose Zeit zwischen Infektion (Eintritt d.
  Erregers in den Organismus) und dem Auftreten erster
  Symptome.
  \item[\T{Epidemie}]\index{Epidemie}
  örtlich begrenztes, zeitlich vermehrtes Auftreten einer 
  Infektionserkrankung (z.\,B. Grippe, Salmonellen).
  \item[\T{Pandemie}]\index{Pandemie}
  Epidemie die sich über Länder u. Kontinente ausbreitet
  (z.\,B. Pest im Mittelalter).
  \item[\T{Endemie}]\index{Endemie}
  Auf kleinere Gebiete beschränktes, zeitlich jedoch nicht
  begrenztes Auftreten einer Infektionserkrankung  (z.\,B.
  FSME,
  Malaria, Lepra).
  \item[\T{Morbidität}]\index{Morbidität}
  Anteil der an einer Erkrankung \EE{Erkrankten} einer bestimmten
  \E{Bevölkerungsgruppe} (z.\,B. 5 pro 100\,000 Einwohner).
  \item[\T{Mortalität}]\index{Mortalität}
  Anteil der an einer Erkrankung \EE{Verstorbenen} einer
  bestimmten \E{Bevölkerungsgruppe} (z.\,B. 5 pro 100\,000
  Einwohner).
  \item[\T{Letalität}]\index{Letalität}
  Anteil der  an einer Erkrankung \EE{Verstorbenen} bezogen
  auf die \E{Gesamtzahl der Erkrankten} (\Speech{\enquote{1
  von 5 Erkrankten stirbt}}).
  \item[\T{Entzündung}] \prettyref{topic:entzuendung};
  \begin{inparaenum}
    \item Rötung,
    \item Schwellung,
    \item Überwärmung,
    \item Schmerzen,
    \item Funktionsverlust.
  \end{inparaenum}
  \item[\T{Prophylaxe}]\index{Prophylaxe} 
  Vorbeugende Maßnahme(n).
\end{labeling}

%\end{multicols}





%%%%%%%%%%%%%%%%%%%%%%%%%%%%%%%%%%%%%%%%%%%%%%%%%%%%%%%%%%%
%%%%%%%%%%%%%%%%%%%%%%%%%%%%%%%%%%%%%%%%%%%%%%%%%%%%%%%%%%%
%%%%%%%%%%%%%%%%%%%%%%%%%%%%%%%%%%%%%%%%%%%%%%%%%%%%%%%%%%%
\subsection{Die wichtigsten Erreger}
\label{topic:infektion-erreger}

\begin{table}[H]\FontTableBase
\Captionof{table}{Übersicht der Erregerarten.}

\begin{tabulary}{\linewidth}{LL}\addlinespace\toprule\addlinespace
\TabHeading{Erreger}
&
\TabHeading{Kurzbeschreibung}
\\\addlinespace\midrule\addlinespace
\TabSub{Bakterien}
&
klein, im Mikroskop meist sichtbar;
        {\textquotedblleft}Normale{\textquotedblright} Zellen
\\\addlinespace
\TabSub{Viren}
&
sehr klein, im Mikroskop i.\,d.\,R. nicht
        sichtbar, keine echten Zellen, müssen Zellen befallen
\\\addlinespace
\TabSub{Pilze}
&

\\\addlinespace
\TabSub{Parasiten, Tiere}
&
Würmer, Maden, \ldots
\\\addlinespace
\TabSub{(Toxine)}
&
Giftstoffe, werden z.\,B. von Bakterien produziert.
\\\addlinespace\bottomrule
\end{tabulary}
\end{table}

\subsubsection{Bakterien}
\index{Bakterien}

\TextSlide{\TxBeschreibung}{%
Bakterien sind vollständige,
'echte' Zellen, welche leben
und aktiv etwas tun können. Dazu gehört z.\,B. die Produktion
von diversen -- z.\,T. schädlichen -- Substanzen (Toxine, ...) und Gasen etc.
\Z{Sie können die unterschiedlichsten Formen (Kugeln, Stäbchen, mit Geißeln \ldots ) haben.}

Bakterien sind nicht grundsätzlich schlecht. Viele können
für den Menschen gut, neutral oder manchmal auch schädlich sein.
Dies hängt von verschiedenen Faktoren ab, s.\,u.
}{
\begin{itemize}
  \item Vollständige, \protect\enquote{echte}
  Zellen\,/\,Lebewesen
  \item Nicht grundsätzlich schlecht
\end{itemize}
}{}{}{}




\TextSlide{Bakterien und Pathogenität}
{%
Unter \E{Pathogenität} versteht man die Eigenschaft eines
Bakteriums, eine Erkrankung oder sonstigen Schaden
auszulösen. Die Schädlichkeit von Bakterien hängt von
verschiedenen Faktoren ab. Grundsätzlich entscheidend ist
die \EE{Bakterienart}: Je nachdem welcher Spezies und
welchem Stamm das Bakterium angehört, kann es
unterschiedliche Wirkungen haben. Der \EE{Ort der
Besiedelung} ist ebenso entscheidend. Viele Bakterien haben
angestammte Besiedlungsorte, wo sie natürlich vorkommen oder
sogar vorkommen müssen, damit eine gewünschte Funktion
(z.\,B. die Verdauung) funktioniert. Damit eine Bakterienart
nicht ungehemmt wuchern kann, ist der
\enquote{\EE{Konkurrenzkampf der Bakterien}} notwendig. Wenn
sich verschiedene Bakterienarten um den gleichen Raum und
die gleichen Nährstoffe 'streiten' müssen, können einzelne
Stämme keine derart schädigende Wirkung entfalten, als wie
wenn sie ungestört und ungehemmt wachsen könnten.

Die \EE{Abwehrkräfte des Patienten} (Immunstatus) sind
gerade im medizinischen Umfeld ein wichtiger Faktor.
Normalerweise kommt das Immunsystem des Menschen gut mit den
natürlich vorkommenden Keimen zurande. Wenn das
\E{Immunsystem} jedoch \E{geschädigt oder geschwächt} ist,
können  \E{auch sonst harmlose Keime} zu einer großen Gefahr
werden! Besonders betroffen sind Patienten mit einer
\E{chemotherapeutischen Behandlung} oder
Leukämie\footnote{\E{Leukämie}:
Bösartige Erkrankung des blutbildenden Systems, es kommt
u.\,a. zu einer unzureichenden Bildung von für das
Immunsystem wichtigen weißen Blutkörperchen.}, sowie
Patienten mit \E{Immunschwäche}-Erkrankungen (z.\,B.
\Abk{AIDS}, \prettyref{topic:aids}). Auch \E{Kinder},
\E{ältere Personen} oder Patienten in schlechter
körperlicher Verfassung haben oft ein schwächeres
Immunsystem. Auch sonst gesunde Menschen können aufgrund von
Stress oder schlechter Ernährung anfälliger für Krankheiten
sein.
}{
\begin{itemize}
  \item Bakterienart
  \item Keimzahl
  \item Ort der Besiedelung
  \item Konkurrenzkampf
  \item Abwehrkräfte des Patienten
  \begin{itemize}
    \item Kind, Chemotherapie, AIDS, ...
    \item Körperliche Verfassung
    \item Alter, Ernährung, Stress, ...
  \end{itemize}
\end{itemize}
}
{images/AASdev20090228-img96.jpg}
{Für manche Bakterien zahlt man
    freiwillig viel Geld, hier am Beispiel des
    Lactobacillus casei gezeigt.}
{}



\Fazit{Daraus folgt: Manche Keime machen immer krank, manche
nur unter Umständen.
  \EE{Opportunisten} machen nur krank wenn der Wirt
  geschwächt ist. Das sind viele Patienten im Rettungs- und
  Krankentransport!
}


\TextSlide{Bakterien, Antibiotika und Resistenzen}{%
Antibiotika sind Medikamente, die Bakterien abtöten
oder deren Vermehrung hemmen. Sie haben i.\,d.\,R. ein bestimmtes Spektrum, dh.
wirken nur auf bestimmte Keime.

Ein großes Problem ist die zunehmende Bildung von
\EE{Resistenzen} von Bakterien gegen Antibiotika, welche
diese Medikamente wirkungslos machen. Man spricht auch von
der Bildung von \E{multiresistenten Keimen}. Diese werden
unter \prettyref{topic:multiresistente-keime} näher
besprochen.
}{
\begin{itemize}
  \item Bakterien können Resistenzen gegen Antibiotika bilden. 
\end{itemize}
}{}{}{}






\subsubsection{Viren}

\TextSlide{\TxBeschreibung}{%
\T[Virus]{Viren} sind kleiner als Bakterien und sind keine
vollständigen Zellen.\ZFootnote{Sie sind so klein dass sie
im Lichtmikroskop nicht sichtbar sind.} \Z{Sie brauchen
Wirtszellen um sich vermehren zu können. Sie schleusen ihr
eigenes genetisches Programm (Erbinformationen) in die
fremde Zelle ein, diese produziert dann neue Viren.
Die missbrauchte Zelle verliert i.\,d.\,R. ihr eigentliche
Funktion oder wird komplett zerstört.}

Virale Infektionen sind sehr häufig. Die meisten sog.
\EE{\enquote{Verkühlungskrankheiten}}, sowie andere bekannte
Krankheiten wie z.\,B. die \EE{Grippe}, \E{Feuchtblattern}
und viele \E{Hepatitis}-Formen, werden durch Viren
verursacht.
}{
\begin{itemize}
  \item Keine vollständigen Zellen
  \item Benötigen Wirtszellen zur Vermehrung
  \item Virale Erkrankungen sind häufig:
  \begin{itemize}
    \item Grippe
    \item Feuchtblattern
    \item Hepatitis-Formen
    \item $\cdots$
  \end{itemize}
\end{itemize}
}{}{}{}




\TextSlide{Nachweis}{%
Die Tests haben üblicherweise ein \E{diagnostisches
Fenster}, in dem eine Infektion (noch) nicht nachweisbar
ist.
\ZFootnote{Viren lassen sich nicht auf Nährböden anzüchten.
Man verwendet entweder indirekte Antigen- oder
Antikörpernachweistests oder die Polymerase-Kettenreaktion
(PCR) zum direkten Nachweis der Virus-DNA/RNA. Es dauert
jedoch einige Zeit bis nach einer Infektion ausreichend
entsprechendes Material gebildet wird um den Test positiv
werden zu lassen.}
}{
\begin{itemize}
  \item Diagnostisches Fenster: Test wird erst einige Zeit nach der Infektion
  aussagekräftig.
\end{itemize}
}{}{}{} 


%%%%%%%%%%%%%%%%%%%%%%%%%%%%%%%%%%%%%%%%%%%%%%%%%%%%%%%%%%%
%%%%%%%%%%%%%%%%%%%%%%%%%%%%%%%%%%%%%%%%%%%%%%%%%%%%%%%%%%%
%%%%%%%%%%%%%%%%%%%%%%%%%%%%%%%%%%%%%%%%%%%%%%%%%%%%%%%%%%%
\subsection{Die wichtigsten Übertragungswege}
\label{topic:infektion-uebertragungswege}



Krankheitserreger können auf unterschiedliche Wege in den
Körper gelangen und eine Infektion verursachen:

\begin{multicols}{2}
  \begin{itemize}
    \item Fäkal-oral (Schmierinfektion, z.\,B. Badesee)
    \item Aerogen (Tröpfcheninfektion, Aerosol)
    \item Genital (Geschlechtsverkehr)
    \item Über/durch die Haut
    \item Diaplazentär (während Schwangerschaft)
    \item Während der Geburt
    \item Eintrittspforten: alles was irgendwie offen ist
    \begin{itemize}
      \item Wunden, Magen-Darm-Trakt, Atmung, Auge,
      Penis, Scheide,
      Plazenta
      \ldots
    \end{itemize}
    \item Fremdkörper sind Wohnhäuser und Autobahnen für Keime
    \begin{itemize}
      \item Harnkatheter, venöse Zugänge, Tubus, künstlicher
      Mageneingang, Dialysekatheter {\dots}
    \end{itemize}
  \end{itemize}
\end{multicols}




%%%%%%%%%%%%%%%%%%%%%%%%%%%%%%%%%%%%%%%%%%%%%%%%%%%%%%%%%%%
%%%%%%%%%%%%%%%%%%%%%%%%%%%%%%%%%%%%%%%%%%%%%%%%%%%%%%%%%%%
%%%%%%%%%%%%%%%%%%%%%%%%%%%%%%%%%%%%%%%%%%%%%%%%%%%%%%%%%%%
\subsection{Kreuzinfektion}
\index{Kreuzinfektion}
\label{topic:kreuzinfektion}

\TextSlide{\TxBeschreibung}
{%
Unter einer \T{Kreuzinfektion} versteht man die
Verschleppung von Erregern von einem Patienten auf einen
(unbeteiligten) anderen, wobei das Personal oder Gegenstände
als \enquote{Vehikel} fungieren. \E{Hygienisches Arbeiten}
und \E{sachgerechte Desinfektionsmaßnahmen} sind wichtige
Mittel zur \E{Vermeidung von Kreuzinfektionen}.

\Fazit{Kreuzinfektion: \EE{Patient 1
${\rightarrow}$ \emph{Personal} ${\rightarrow}$ Patient
2}}
}{
Patient 1 ${\rightarrow}$ \emph{Personal} ${\rightarrow}$
Patient 2 }{}{}{}


\subsection{Angepasste Hygiene}

Grundsätzlich müssen die Hygienevorgaben und -maßnahmen dem
jeweiligen Anforderungen und Gegebenheiten
angepasst werden: Bei einem zu Hause lebenden
Krankentransportpatienten bestehen zum Beispiel andere
Risiken als bei einem beatmeten Intensivpatienten.


\opt{STATUS-EXPERIMENTAL}{
\Layout{57}{}{}{}{
\begin{tabularx}{\linewidth}{XXXX}

&
Normalbürger
&
Krankentransportpatient
&
Intensivpatient
\\\addlinespace
\TabSub{Problembereiche}
&

&

&
Beatmung

Antibiotikatherapie

Wunden

Invasivitäten: Sonden (Magen, PEG, Hindruckmessung,
{\ldots}), Katheter, Gefäßzugänge (zentral und peripher), \ldots 
\\\addlinespace
\TabSub{Risken für Patient}
&

&
Wunden

Kreuzinfektionen

Resistente Keime

Grunderkrankung; evtl. Immunschwäche, evtl. schlechter
Allgemeinzustand $\pm$ höhere Infektanfälligkeit 
&
Resistente Keime

Verkeimtes Beatmungszubehör

Wunden

Schlechter Allgemeinzustand, höhere Infektanfälligkeit

Immunschwäche
\\\addlinespace
\TabSub{Risken durch Pateint}
&

&
Infektiöse Erkrankungen

Evtl. Besiedelung durch resistente
Keime 
&
Infektiöse Erkrankungen

Hohe Wahrscheinlichkeit für Besiedelung durch resistente
Keime 
\\\addlinespace
\end{tabularx}
}
{Beispiel für Unterschiedliche Hygiene-Risikoprofile}
{}
}%STATUS-EXPERIMENTAL


% \section{Wichtige Krankheitsbilder}
%\marginlabel{Krankheitsbilder}
\clearpage
\section{Multiresistente Keime}
\label{topic:multiresistente-keime}

\TextSlide{Linktip}{%
\cite{mrsanet} \emph{EUREGIO MRSA-net}
\url{http://www.mrsa-net.org/}
}{}{}{}{}

\Layout{60}{Geschichte und Bedeutung der Antibiotika}
{%
Vor der Entdeckung der Antibiotika im 20. Jahrhundert waren
bakterielle Infektionen sehr gefürchtet, viele Menschen sind
an -- nach heutiger Sicht simplen -- Infektionen gestorben.
Wer es im Ersten Weltkrieg lebendig in ein Lazarett geschafft
hatte, verstarb dort im Verlauf  oft an den Folgen des
\E{Wundbrandes}. Erkrankungen wie \E{Scharlach} oder
\enquote{Blutvergiftungen} waren häufig ein Todesurteil.

Im Jahr 1928 verschimmelte dem englischen Wissenschaftler
Alexander Fleming versehentlich eine Nährplatte mit
Bakterienkulturen. Er bemerkte, dass in der Umgebung des
Pilzes die Bakterien nicht wachsen konnten: das
\E{Penicillin} war entdeckt.\ZFootnote{Zuvor hatten bereits
deutsche Wissenschaftler eine andere Antibiotika-Art
entwickelt, die Sulfonamide.} Der erste Patient der mit
Penicillin behandelt wurde war ein Londoner Polizist. Er
hatte sich beim Rasieren verletzt und eine Blutvergiftung
zugezogen -- er galt als unrettbar.~\ZFootnote{Leider wusste
man damals noch nicht, wieviel Penicillin ein Mensch für
eine erfolgreiche Behandlung benötigte. Der Polizist
verstarb -- nach anfänglicher deutlicher Besserung -- als
das experimentell hergestellte Penicillin aufgebraucht war.}

In weiterer Folge wurde eine Vielzahl weiterer Antibiotika
entwickelt und man dachte, man hätte die bakteriellen
Infektionen besiegt. Doch man irrte sich gewaltig {\ldots}
}
{}
{./images/hygiene/Faroe_stamp_079_europe_fleming.jpg}
{\Z{Alexander Fleming.}}
{PD, Postverk Føroya -
Philatelic Office}



\TextSlide{Resistenzen}
{%
Bakterien können Resistenzen gegen Antibiotika entwickeln,
welche dann nicht mehr oder nur noch sehr abgeschwächt
helfen. Als behandelnder Arzt steht man dann hilflos vor dem
Patienten.
Es muss daher das Ziel aller sein die weitere Ausbreitung von
resistenten Keimen zu verhindern. Dem Rettungs- und
Krankentransport kommt dabei eine \E{Schlüsselrolle} zu, da
es hier leicht zu \EE{Kreuzinfektionen} kommen kann.

Gesundheitseinrichtungen wie Spitäler oder Pflegeheime sind
eine ideale \enquote{Zuchtanstalt} und Brutstätte für Resistenzen:
Es werden viele Antibiotika verwendet, viele Patienten liegen
räumlich nahe beieinander. Es finden sich daher dort
\E{gehäuft} resistente Keime, umgangssprachlich spricht man
in der Fachwelt auch von
\index{Spitalskeime}\T{Spitalskeimen}.
Resistente Keime verursachen komplizierte Infektionen, eine
Erhöhung der Komplikations- und Sterblichkeitsrate, in Folge
eine Verlängerung des Krankenhausaufenthaltes und nicht
zuletzt enorme zusätzliche -- vermeidbare -- Kosten
\cite{Kucera2010}.
}{
\begin{itemize}
   \item Bakterien können Resistenzen entwickeln!
   \item Antibiotika helfen dann nicht mehr!
   \item Gesundheitseinrichtungen besonders betroffen
\end{itemize}
}
{}
{}
{}

\Fazit{Resistenzen sind ein enorm großes Problem in der
heutigen Medizin!}

\Text{Im Krankentransport gibt es Problemkeime}
{%
Da im Krankentransport natürlich häufig Patienten aus
Gesundheitseinrichtungen (oder zwischen
Gesundheitseinrichtungen) transportiert werden, ist man auch
hier mit dem Problem der resistenten Keime besonders
konfrontiert. Deshalb sind spezielle Maßnahmen erforderlich.

Der wohl bekannteste Problemkeim ist der \T{MRSA}
(\Z{Multiresistenter Staphylococcus aureus, siehe unten)}.
Das liegt schlicht daran, dass er der erste war. In den
letzten Jahren ist die Resistenzentwicklung aber auch bei
vielen anderen Keimen zu beobachten, \Z{Beispiele dafür sind
die \Abk{ESBL}{\textquotesingle}s und
\Abk{VRE}{\textquotesingle}s, sowie der VRSA (Extended
Spectrum Beta-Laktamase-Bildner, Vancomycin-resistente
Enterokokken, Vancomycin-resistenter Staph. aureus). Auch bei
der Tuberkulose treten mittlerweile resistente Keime auf:
\E{2008} wurden die ersten Fälle registriert,  bei denen der
Keim auf mind. \E{4} Antibiotika resistent war\opt{NIVMED}{
(\E{XDR-TBC}. Resistenzen gegen: Isoniazid, Rifampizin,
Fluochinolone, sowie zumindest eines der folgenden
Antibiotika: Amikacin, Capreomycin,
Kanamycin)}}\footnote{Quelle: AGES}. \E{Das Problem der
multiresistenten Keime wird immer größer.}

\TODO{GABS}{yyyy-mm-dd}{LIT, HYG: MDR/XDR-TB Quelle bei AGES}{}

}
{
\begin{itemize}
  \item Krankentransport: Patienten aus
  Gesundheitseinrichtungen
  \item Problemkeime häufig
  \item MRSA häufig
\end{itemize}
}
{}
{}
{}










\Fazit{Resistenzen machen Antibiotika zunehmend wirkungslos.
  -- Bald sind wir wieder dort, wo wir 1927, vor der
  Entdeckung des Penicillins, waren: Praktisch machtlos gegen
  bakterielle Infektionen.}
\Fazit{Je wirkungsloser die Antibiotika werden, desto
  wichtiger ist die Hygiene und die Vermeidung von
  (Kreuz-)Infektionen.}
\Fazit{Auch im Rettungs- und Krankentransport kommen
  gehäuft resistente Keime vor! Ein Haltegriff in der
  U-Bahn ist daher nicht mit dem Haltegriff im
  Krankentransportwagen
  vergleichbar.}




\subsection{MRSA -- Multi-Resistenter Staphylococcus
aureus} 

\Z{Ursprünglich: Methicillin-Resistenter Staphylococcus Aureus}


\Layout{60}{Staphylococcus aureus}
{%
Der \emph{Staphylococcus aureus} ist ein
\enquote{Allerweltskeim}, bis zu 50\PC* der Bevölkerung sind
besiedelt. Vor allem ist die vordere Nasenhöhle besiedelt,
aber auch die Haut (Achselhöhle, Analregion) und die
Schleimhäute können ihn beherbergen. Auch im Spital ist er
häufig, aufgrund des hohen Antibiotikaverbrauchs ist er dort
aber oft schon Antibiotika-resistent: Man spricht dann vom
\T{MRSA} (Multi-resistenter Staphylococcus aureus).

\vspace{1cm}

}{
\begin{itemize}
    \item Staph. aureus = \enquote{Allerweltskeim}
    \item aber resistente Form:
    \T{MRSA}
%     \begin{itemize}
%         \item {\textquotedblleft}Spitalsinfektionen{\textquotedblright}
%         \item schwer bis gar nicht behandelbar
%         \item Wunden, die nicht abheilen wollen, Sepsis, Tod ...
%     \end{itemize}
\end{itemize}
}
{./images/hygiene/11159_lores_VRSA.jpg}
{\AuthNotesPicLegalOk{PD}{}Ein
resistenter Staph. aureus (\Abk{VRSA}) im Elektronenmikroskop,
koloriert.}
{Public Health Image Library, ID\#: 11159, Public Domain.}



\begin{PictureSeriesWide}{}{Multiresistente Keime in der
Praxis. Auf Nährmedien können Bakterien zu diagnostischen
Zwecken gezüchtet werden. Ebenso kann man die Wirkung von
Antibiotika auf den Keim untersuchen.}
\PictureSeriesRowWide{3}
{./images/gabriel-christine-aass/staph-aureus.jpg}
{Auf einem Nährboden wurde eine Bakterienkultur
angelegt. Jedes der sechs runden Blättchen enthält ein Antibiotikum. Man kann
sehen dass in der Umgebung der Blättchen keine Bakterien wachsen.
\SourcePicture{Gabriel/KARPAT}}
{./images/gabriel-christine-aass/staph-aureus-mrsa.jpg}
{Diese Bakterien ignorieren 4 von 6
Antibiotika-Blättchen, d.\,h. die entsprechenden Antibiotika wirken
nicht gegen diesen Keim. Er ist \EE{multiresistent}.
\SourcePicture{Gabriel/KARPAT}}
{./images/gabriel-christine-aass/amputat-us1.jpg}
{Oft die Folge einer infizierten, nicht abheilenden
Wunde: Eine Amputation des Unterschenkels.
\SourcePicture{Gabriel/KARPAT}}
{empty}
{}
\end{PictureSeriesWide}



% M %%%%%%%%%%%%%%%%%%%%%%%%%%%%%%%%%%%%%%%%%%%%%%%%%%% M %
% Multiresistente Keime
\ProcedurePrintDefault{IU81001C}



\clearpage
\section{Reinigung und hygienische Wiederaufbereitung}



\subsection{Allgemeine Bemerkungen zur Wiederaufbereitung}

Reinigung, Desinfektion und Sterilisation von
Medizinprodukten in Einrichtungen des Gesundheitswesens sind
mit Geräten und geeigneten validierten Verfahren so
durchzuführen, dass der Erfolg dieser Verfahren
nachvollziehbar gewährleistet ist und die Sicherheit und
Gesundheit von Patienten, Anwendern oder Dritten nicht
gefährdet wird. \E{(Desinfektions-)Mittel}, welche zur
Wiederaufbereitung von Medizinprodukten eingesetzt werden,
gelten wiederum als \E{Medizinprodukte}\footnote{{\S}\,2 Abs
2 Medizinproduktegesetz}; es gelten daher die gleichen
Unterweisungspflichten.


\subsection{Techniken zur Wiederaufbereitung}

Man unterscheidet grundsätzlich zwischen \E{Reinigung}, \E{Desinfektion} und 
\E{Sterilisation}.

\subsubsection{Reinigung}

\TextSlide{\TxBeschreibung: Reinigung}
{%
Unter \T{Reinigung} versteht man das Entfernen von groben --
oft sichtbaren -- Verunreinigungen mit vorrangig
mechanischen Mitteln. Sie muss möglichst \EE{rasch nach der
Kontamination begonnen} werden, da es sonst zum Eintrocknen
der Verunreinigungen kommen kann. Gereinigte Gegenstände
können noch infektiös sein!
Eine (Vor-)Reinigung ist Voraussetzung für eine
weiterführende Desinfektion oder Sterilisation.
}{
\begin{itemize}
  \item Entfernung von groben Verunreinigungen
  \item Beginn rasch nach Kontamination
  \item Voraussetzung für spätere Desinfektion od.
  Sterilisation
\end{itemize}
}{}{}{}



\subsubsection{Desinfektion}
\index{Desinfektion}
\label{topic:desinfektion-grundzuege}

\TextSlide{\TxBeschreibung: Desinfektion}
{%
Unter dem Begriff \T{Desinfektion} versteht man Maßnahmen
die der \EE{Keimreduktion} dienen, um einen Gegenstand so
aufzubereiten, dass er \underline{unter Normalbedingungen}
nicht infektiös ist. Der Gegenstand ist dabei jedoch
\E{nicht} keimfrei. Die Desinfektionsmaßnahmen töten
i.\,d.\,R. nur bestimmte Keime ab. Weit verbreitet ist die
Desinfektion mittels chemischer Substanzen
(Desinfektionslösungen, Einlegebäder, \ldots). Dabei muss
unbedingt die korrekte Konzentration der Lösung und die
vorgeschriebene Einwirkzeit beachtet werden.
Eine Desinfektion ist i.\,d.\,R. für Gegenstände
ausreichend, die \EE{nur mit intakter Haut} in Berührung
kommen.
Voraussetzung für eine Desinfektion ist eine
\EE{vorhergehende Reinigung} des Gegenstandes.
}{
\begin{itemize}
  \item Keimreduktion
  \item Z.\,B. mittels chem. Substanzen
  \item Vorher: Reinigung!
\end{itemize}
}{}{}{}

\Fazit{Keine Desinfektion ohne vorhergehender Reinigung!}

\clearpage 
\TextSlide{Wirkung von Desinfektionsmitteln}
{%
Desinfektionsmittel haben grundsätzlich eine keimabtötende
Wirkung. Diese ist jedoch z.\,T. sehr stark von der Art des
Desinfektionsmittels, der Art der Erreger und anderen
Faktoren abhängig.

\subparagraph{Einsatzbereich} 
Jedes Desinfektionsmittel hat einen klar umrissenen
Einsatzbereich (Flächendesinfektionsmittel,
Instrumentendesinfektionsmittel, Desinfektionsmittel für
Einlegebäder, Händedesinfektionsmittel,
Hautdesinfektionsmittel, \ldots{}). Ein Desinfektionsmittel
darf nur im dafür vorgesehenen Bereich verwendet werden, da
es sonst wirkungslos oder schädlich sein kann.\footnote{So
darf man zum Beispiel nicht vor einer Venenpunktion die Haut
mittels eines Händedesinfektionsmittel desinfizieren.
Händedesinfektionsmittel beinhalten rückfettende Substanzen
um hautverträglicher zu sein. Gelangen diese durch den Stich
in das Blutgefäß kann es zu Venenreizungen kommen.}
Verschiedene Flächen können auch unterschiedliche
Desinfektionsmittel benötigen: So dürfen Acryl-Oberflächen
keinesfalls mit Alkohol-hältigen Mitteln desinfiziert werden,
da sonst das Material angegriffen und stark geschädigt wird.
Es muss für jede Oberfläche ein passendes Mittel verwendet
werden.

Die \E{Pflegehinweise des Herstellers} des zu reinigenden
Produktes sind zu beachten! Hinweise zur Reinigung und
hygienischen Wiederaufbereitung eines Medizinproduktes
müssen Bestandteil  jeder \E{Einschulung} sein und müssen
auch in der \E{Gebrauchsanweisung} zu finden sein.
In Fahrzeugen dürfen i.\,d.\,R. aufgrund der
Acryloberflächen nur Flächendesinfektionsmittel ohne Alkohol
verwendet werden (z.\,B. \I{Acryl-Des{\texttrademark}},
\I{Acrylan{\texttrademark}})!


\subparagraph{Einwirkzeit} Der wichtigste Faktor ist die
Einwirkzeit: Erst nach einer gewissen Zeit beginnt das Mittel
zu wirken und tötet tatsächlich Keime ab. Hält man die
notwendige Zeit nicht ein, so ist das Desinfektionsmittel
nahezu \E{wirkungslos}. Die erforderliche Zeit ist den
jeweiligen Herstellerangaben zu entnehmen und ist für jedes
Produkt unterschiedlich!

\subparagraph{Konzentration} Manche Desinfektionslösungen
müssen vom Anwender selber durch \E{Verdünnung eines
Konzentrates} hergestellt werden.  Die vorgeschriebene
Verdünnung muss exakt eingehalten werden, da es sonst zu
einem Wirkungsverlust kommt.


\subparagraph{Unsachgemäße Anwendung} Die unsachgemäße
Herstellung, Anwendung oder Pflege von Desinfektionsmitteln
und -lösungen kann \E{katastrophale Folgen} haben. Im besten
Fall tritt nur ein Wirkungsverlust ein, es kann jedoch auch
dazu kommen dass sich in der Desinfektionslösung
\EE{Keime sogar vermehren!} 

}{
\TabSub{Einsatzbereich}

\begin{itemize}
  \item Flächen, Instrumente, Hände \ldots{}
  \item Jedes Mittel darf nur für sein jeweiliges
  Einsatzgebiet verwendet  werden
\end{itemize}

\TabSub{Einwirkzeit}

\begin{itemize}
  \item Unbedingt beachten, sonst Wirkungsverlust
\end{itemize}

\TabSub{Konzentration}

\begin{itemize}
  \item Unbedingt beachten, sonst Wirkungsverlust
\end{itemize}

\TabSub{Unsachgemäße Anwendung, Herstellung, Pflege}

\begin{itemize}
  \item Gefährlich
  \item Im Extremfall können Keime \E{im Desinfektionsmittel} wachsen
\end{itemize}
}{}{}{}



\Cave{In Fahrzeugen dürfen i.\,d.\,R. keine alkoholischen
Desinfektionsmittel verwendet werden, da diese die
Acryloberflächen zerstören können!}


\Cave{Die Verdünnungsangabe \enquote{\Ca{1:10}} steht für \enquote{1~Teil
Desinfektionsmittel und
\Ca{\underline{9}~Teile Wasser!}}} 

\Fazit{Desinfektionsmittel müssen richtig hergestellt,
gepflegt und angewendet werden, sonst können sie mehr
Schaden als Nutzen bringen! }

\TextSlide{Desinfektion mittels Einlegebad}
{%
\E{Es handelt sich hierbei um eine schlechte
Möglichkeit zur Desinfektion, da der Erfolg der Desinfektion
von vielen variablen Faktoren abhängt und eine
zufriedenstellende Leistung kaum sichergestellt werden kann
(nicht validierbares Verfahren). Sie ist daher nur dann zulässig,
wenn eine geringe Infektionsgefahr besteht und keine andere
Methode sinnvoll durchführbar ist.
Für Instrumente ist eine Desinfektion mittels Einlegebad
ist nur zulässig, wenn diese ausschließlich in Kontakt mit
intakter Haut kommen.}

Bei der Desinfektion
mittels Einlegebad wird in einer geeigneten Wanne eine Desinfektionslösung
zubereitet\footnote{Typischer Weise wird ein Konzentrat
entsprechend verdünnt.} und Gegenstände darin über einen
längeren Zeitraum (mehrere Stunden) eingelegt.
Dabei ist es wichtig, dass die gesamte Oberfläche benetzt
ist, es muss daher \EE{luftblasenfrei} eingelegt werden.
Schwimmende Teile werden mittels eines \EE{Niederhaltesiebes}
unter den Flüssigkeitsspiegel gedrückt. Zerlegbare
Gegenstände müssen \EE{soweit als möglich zerlegt} werden;
Scheren, Klemmen und ähnliches werden geöffnet eingelegt. Die
\EE{Einwirkzeit} muss strikt eingehalten werden. Nach der
Einlegedesinfektion ist i.\,d.\,R. eine \EE{Endreinigung}
durchzuführen, um Rückstände des Desinfektionsmittels zu
beseitigen.



\subparagraph{Haltbarkeit}
Die hergestellte Lösung ist nur \E{begrenzt haltbar} und muss
regelmäßig getauscht werden. Ein Tausch ist auch unbedingt
notwendig, wenn die Lösung sichtlich getrübt oder verschmutzt
ist!

\Cave{In unsachgemäß gepflegten Desinfektionsmitteln können Keime wachsen!}

}{
\begin{itemize}
  \item Luftblasenfrei
  \item Vollständig untertauchen, evtl. mit Niederhaltesieb
  \item Einwirkzeit beachten! 
  \item Endreinigung
  \item Desinfektionslösung regelmäßig, bzw. auch bei Verunreinigung, wechseln
\end{itemize}
}{}{}{}

\TextSlide{Flächen- und Wischdesinfektion}
{%
\index{Flächendesinfektionsmittel}
\index{Acryl-Des{\texttrademark}}
\index{Acrylan{\texttrademark}}
Eine Flächendesinfektion wird mit dafür vorgesehenen
\EE{Flächendesinfektionsmitteln} durchgeführt. Dabei ist
darauf zu achten, dass für \EE{jede Oberfläche das richtige
Mittel} verwendet wird. So zerstören z.\,B.
Alkohol-haltige Desinfektionsmittel Acryl-Oberflächen, wie
sie oft in Fahrzeugen oder bei medizinischen Geräten
vorzufinden sind (s.\,o.). Es werden in so einem Fall
Alkohol-freie Mittel (z.\,B. Acryl-Des{\texttrademark},
Acrylan{\texttrademark}) angewendet.

Auch hier ist die vom Hersteller vorgegebene
\EE{Einwirkzeit} zu beachten, die Fläche muss während dieser
Zeit \EE{feucht} gehalten werden!

Um den Schmutz nicht zu verteilen wird \EE{von außen
in Richtung einer Schmutzquelle gewischt}. Es dürfen nur
Einweg-Putztücher verwendet werden, diese sind regelmäßig zu
wechseln. Die Verwendung von Putzlappen zur
Wischdesinfektion ist \E{nicht zulässig}!


\Cave{In den Fahrzeugen dürfen keine alkoholischen
Desinfektionsmittel verwendet werden, da diese die
Acryloberflächen zerstören können!}
}{
\begin{itemize}
  \item Mit Flächendesinfektionsmittel
  \item Einwirkzeit beachten! Fläche muss während der Zeit
  \E{feucht} sein!
  \item Material-entsprechende Substanzen!
  \begin{itemize}
    \item z.\,B.: Alkohol zerstört Acryl-Flächen! ${\rightarrow}$
    Alk.-freie Mittel
  \end{itemize}
  \item \EE{Von außen \E{in Richtung Schmutzquelle}
  wischen}
  \item Einweg-Tücher, regelmäßig wechseln
\end{itemize}
}{}{}{}

\subsubsection{Sterilisation}
\label{topic:sterilisation-grundzuege}
\TextSlide{\TxBeschreibung: Sterilisation}
{%
Das Ziel der \T{Sterilisation} ist die Abtötung
\E{sämtlicher} Krankheitserreger, d.\,h. \EE{Keimfreiheit}
(Sterilität).
Voraussetzung für die Sterilisation ist eine \E{vorherige
Reinigung} und \E{Desinfektion}!

\Z{Das bekannteste Verfahren ist das \E{Autoklavieren},
dabei wird der Gegenstand unter Überdruck mit Hitze und
Wasserdampf behandelt.} Sterilisationsverfahren sind relativ
aufwendig und setzen Wissen und Erfahrung
voraus.\footnote{Sterilgut sollte nur von entsprechend
qualifizierten und zertifizierten Einrichtungen mit
validierten Verfahren aufbereitet werden!}
}{
\begin{itemize}
  \item Keimfreiheit -- Sterilität
  \item Vorher Reinigung und Desinfektion
  \item Aufwendig
  \item \Z{Z.\,B. Autoklavieren.}
\end{itemize}
}{}{}{}


\Layout{57}
{}
{}
{}
{\begin{tabularx}{\linewidth}[H]{XXXX}
\toprule\addlinespace
&
\noindent\TabHeading{Reinigung}
&
\noindent\TabHeading{Desinfektion}
&
\noindent\TabHeading{Sterilisation}
\\\addlinespace\midrule\addlinespace
\noindent\TabSub{Ziel}
&
Mechanische Beseitigung von Schmutz
&
\EE{Keim\underline{reduktion}} 

Richtet sich gegen ausgewählte Keime
&
Abtötung aller Mikroorganismen 
\\\addlinespace
\noindent\TabSub{Ergebnis}
&
Sauber
&
Keimreduziert -- Nicht keimfrei!
&
\EE{Keimfrei} (steril)
\\\addlinespace
\noindent\TabSub{Arten}
&
Abspülen

Durchspülen

Scheuerreinigung

Maschinelle Reinigung

\ldots
&
Wischdesinfektion

Flächendesinfektion

Einlegebad

\ldots
&
Dampfsterilisation (Überdruck)

Heißluftsterilisation
\\\addlinespace
\noindent\TabSub{Anwendungsbereich}
&
Vorreinigung vor Desinfektion oder Sterilisation

Nicht-medizinische Gegenstände oder Flächen
&
Gegenstände, die ausschließlich mit intakter Haut in
Berührung kommen 
&
Gegenstände, die mit Wunden in Berührung kommen bzw. in
Körperhöhlen eingeführt werden 

Verbandsstoffe

\ldots
\\\addlinespace\bottomrule
\end{tabularx}}
{Gegenüberstellung Reinigung -- Desinfektion --
Sterilisation}
{}



\subsection{Praktisches}

\TextSlide{Hygieneplan - Hygienebeauftragter}
{%
\index{Hygieneplan}
\label{topic:hygieneplan}
Der \T{Hygieneplan} regelt die konkreten Hygienemaßnahmen
zur Erkennung, Verhütung und Bekämpfung von Infektionen bzw.
deren Übertragungsmöglichkeiten. Es wird dort festgelegt
was, wann, wie, und womit (wo) zu geschehen hat. Der
Hygieneplan ist allgemein ausgehängt und hat i.\,d.\,R. die
Bedeutung einer dienstlichen Weisung.

Der  \T{Hygienebeauftragte} (bzw. Hygieneverantwortliche)
ist Ansprechpartner für Belange der Hygiene. Er hat
Maßnahmen zu setzen oder zu veranlassen, die der Erkennung,
Verhütung und Bekämpfung von Infektionen, deren Übertragung
und der Gesunderhaltung dienen. Zur Durchführung dieser
Maßnahmen erstellt er u.\,a. den Hygieneplan. Je nach Art
der Einrichtung hat er eine spezielle
Ausbildung\footnote{Hygienefachkraft, Arzt mit Weiterbildung
für Krankenhaushygiene, Facharzt für Hygiene und
Mikrobiologie, spezielle Fortbildungskurse \ldots{}} bzw.
Kenntnisse im Bereich der Infektionslehre und Hygiene.

% Für Wien: {\S} 9 Abs 2 Z 9, {\S} 19 WRKTG


}
{
\begin{itemize}
  \item Hygieneplan: konkrete Bestimmungen
  \item Hygienebeauftragter: Ansprechpartner mit Zusatzausbildung
\end{itemize}
}{}{}{}


%%%%%%%%%%%%%%%%%%%%%%%%%%%%%%%%%%%%%%%%%%%%%%%%%%%%%%%%%%%
%%%%%%%%%%%%%%%%%%%%%%%%%%%%%%%%%%%%%%%%%%%%%%%%%%%%%%%%%%%
%%%%%%%%%%%%%%%%%%%%%%%%%%%%%%%%%%%%%%%%%%%%%%%%%%%%%%%%%%%
\TextSlide{Persönliche Hygiene}
{\label{topic:persoenliche-hygiene}
Die persönliche Hygiene des Fachpersonals ist ein
wesentlicher Schritt zur Vermeidung von (Kreuz-)Infektionen. 
Schmuck, Uhren, Freundschaftsarmbänder sind perfekte
    Nährböden für Keime (warm, feucht durch Schweiß), und
    verhindern gleichzeitig eine effektive Reinigung. Daher
    müssen diese Dinge vor dem Dienst
    ablegt werden! Weiters müssen bei Dienstantritt
die  Hände gewaschen und anschließend eine hygienische
Händedesinfektion durchgeführt werden.\ZFootnote{Die
Händewaschung dient dem mechanischen Entfernen von Schmutz
und der Verringerung der Anzahl von Bakteriensporen,
welchedurch die alkoholische Händedesinfektion nicht
abgetötet werden können.} Die Nägel und Nagelfalze
müssen optisch sauber sein. Künstliche Fingernägel
oder Nagellack sind nicht statthaft, da sich auf den
künstlichen Oberflächen leicht Keime ansiedeln
können. Lange Haare müssen nach hinten gebunden werden.

Die Dienstkleidung muss gemäß Hygieneplan regelmäßig
gewaschen und natürlich bei Kontamination gewechselt werden.
Wenn mit einer Kontamination gerechnet werden kann, sollen
nach Möglichkeit Einwegschürzen verwendet werden.


}
{\begin{itemize}
    \item Schmuck, Uhren, Freundschaftsarmbänder sind perfekte
    Nährböden 
    \item Kein Nagellack od. künstliche Fingernägel
    \item Haare nach hinten binden
    \item Dienstkleidung regelmäßig waschen! Nach
    Kontamination wechseln!
    \item Plastikschürzen
    \item Bei Dienstantritt
    \begin{compactitem}
        \item Hände waschen
        \item Nagelfalze reinigen
        \item Hygienische Händedesinfektion
        \index{Händedesinfektion}\index{Hygienische Händedesinfektion}
    \end{compactitem}
\end{itemize}}
{}{}{}



\SlideCompensation{2}

\SlideCompensationNoParskip{2}
% \marginpar{\includegraphics[width=4cm]{images/AASdev20090228-img99.jpg}}



% 
% %%%%%%%%%%%%%%%%%%%%%%%%%%%%%%%%%%%%%%%%%%%%%%%%%%%%%%%%%%%
% %%%%%%%%%%%%%%%%%%%%%%%%%%%%%%%%%%%%%%%%%%%%%%%%%%%%%%%%%%%
% %%%%%%%%%%%%%%%%%%%%%%%%%%%%%%%%%%%%%%%%%%%%%%%%%%%%%%%%%%%
% \subsection{Flächen- und Wischdesinfektion}
% 



%%%%%%%%%%%%%%%%%%%%%%%%%%%%%%%%%%%%%%%%%%%%%%%%%%%%%%%%%%%
%%%%%%%%%%%%%%%%%%%%%%%%%%%%%%%%%%%%%%%%%%%%%%%%%%%%%%%%%%%
%%%%%%%%%%%%%%%%%%%%%%%%%%%%%%%%%%%%%%%%%%%%%%%%%%%%%%%%%%%
\TextSlide{Kontaminierte Geräte}
{Unmittelbar nach der Kontamination muss das Eintrocknen von
Schmutzresten verhindert werden. Dazu wird sofort eine
Vorreinigung und -- wenn möglich -- eine Wischdesinfektion
durchgeführt. Schläuche müssen unbedingt mit einer
Reinigungslösung oder evtl. Instrumentendesinfektionsmittel
durchgespült werden.

Sofern kein geeignetes Ersatzmaterial zur Verfügung steht,
ist das Fahrzeug \E{nicht einsatzbereit} und das
kontaminierte Material muss getauscht werden. Auf der
Station wird das Material gemäß Herstellerangaben zerlegt
und noch einmal gespült, und dann der Desinfektion
zugeführt. Üblicherweise erfolgt die Desionfektion mittels
Einlegens in eine Desinfektionswanne. Das Material muss dazu
luftblasenfrei eingelegt und vollständig eingetaucht werden,
ggfs. ist ein Niederhaltesieb zu verwenden. Es muss vermerkt
werden:

\begin{itemize}
  \item Herkunft des Materials (Fahrzeugkennung, \ldots)
  \item Zeitpunkt des Einlegens
  \item Von wem eingelegt
\end{itemize}

Wenn eine sofortige Desinfektion nicht möglich ist (z.\,B.
weil die Desinfektionwanne bereits benutzt wird), muss das
Material in einen dicht verschlossenen Plastiksack verpackt
werden, welcher ebenso beschriftet wird.

\bigskip

\bigskip

\bigskip

\bigskip

\bigskip

\bigskip

\bigskip


}{
\begin{itemize}
  \item Nach Kontamination:
  \begin{compactitem}
    \item \EE{\E{Vorreinigung} vor Ort}: Eintrocknen
    verhindern
    \item Ggfs. \E{Schläuche durchspülen}:     
    \item Evtl. Flächendesinfektion
  \end{compactitem}
  \item Meldung an Leitstelle:
  \emph{{\textquotedblleft}Nicht einsatzbereit, müssen
  putzen{\textquotedblright}}
  \item Kontaminierte Teile gegen Ersatzmaterial tauschen
  \item Zerlegen
  \item Nochmals spülen
  \item Wenn Desinfektionswanne leer ${\rightarrow}$ einlegen
  \begin{compactitem}
    \item \EE{Vollständig zerlegt} laut Anleitung
    \item \EE{Luftblasenfrei!}
    \item Mit Niederhaltesieb unter die Oberfläche drücken
    \item Zeit notieren
    \begin{compactitem}
      \item {\textquotedblleft}was, wann, von wo, von wem
      eingelegt{\textquotedblright}
    \end{compactitem}
    \item sonst nur bereit legen, in Plastiksack
    \begin{compactitem}
      \item \Speech{\enquote*{was, von wo, von wem, warum}}
    \end{compactitem}
  \end{compactitem}
\end{itemize}
}{}{}{}






% \clearpage
%%%%%%%%%%%%%%%%%%%%%%%%%%%%%%%%%%%%%%%%%%%%%%%%%%%%%%%%%%%
%%%%%%%%%%%%%%%%%%%%%%%%%%%%%%%%%%%%%%%%%%%%%%%%%%%%%%%%%%%
%%%%%%%%%%%%%%%%%%%%%%%%%%%%%%%%%%%%%%%%%%%%%%%%%%%%%%%%%%%


\TextSlide{Hygienische Händedesinfektion}
{Die hygienische Händedesinfektion ist eine der wichtigsten
Maßnahmen zur Verhütung von Kreuzinfektionen. Sie ist eine
Standardmaßnahme und wird unter 
besprochen.}
{\CO{[SAN]}
(\FullPrettyRef{topic:hygienische-haendedesinfektion})}{}{}{}









%%%%%%%%%%%%%%%%%%%%%%%%%%%%%%%%%%%%%%%%%%%%%%%%%%%%%%%%%%%
%%%%%%%%%%%%%%%%%%%%%%%%%%%%%%%%%%%%%%%%%%%%%%%%%%%%%%%%%%%
%%%%%%%%%%%%%%%%%%%%%%%%%%%%%%%%%%%%%%%%%%%%%%%%%%%%%%%%%%%
\section{Besondere Verhaltensweisen}

%%%%%%%%%%%%%%%%%%%%%%%%%%%%%%%%%%%%%%%%%%%%%%%%%%%%%%%%%%%
%%%%%%%%%%%%%%%%%%%%%%%%%%%%%%%%%%%%%%%%%%%%%%%%%%%%%%%%%%%
%%%%%%%%%%%%%%%%%%%%%%%%%%%%%%%%%%%%%%%%%%%%%%%%%%%%%%%%%%%




%%%%%%%%%%%%%%%%%%%%%%%%%%%%%%%%%%%%%%%%%%%%%%%%%%%%%%%%%%%
%%%%%%%%%%%%%%%%%%%%%%%%%%%%%%%%%%%%%%%%%%%%%%%%%%%%%%%%%%%
%%%%%%%%%%%%%%%%%%%%%%%%%%%%%%%%%%%%%%%%%%%%%%%%%%%%%%%%%%%
\subsection{Selbstschutz}

Im folgenden seien -- zum Teil als Wiederholung -- die wichtigsten Maßnahmen
zum (hygienischen) Selbstschutz kurz erläutert:


\Text{Schutzkleidung}
{\label{topic:schutzkleidung-hygiene}%
\begin{itemize}
  \item \StepHeading{Untersuchungshandschuhe} sind obligat und
  \EE{immer} bei Patientenkontakt u.\,ä. zu tragen! Auch
  beim Hantieren mit infektiösen oder schädigenden
  Substanzen (darunter fallen auch
  Flächendesinfektionsmittel!) sind Handschuhe zu verwenden.
  Untersuchungshandschuhe bieten keinen 100\PC*igen Schutz,
  bei Risikosituationen empfiehlt sich die Verwendung von
  mind. \E{2 Paar} gleichzeitig.
  Die Handschuhe müssen nach jedem Patienten gewechselt
  werden, nach dem Ablegen ist eine \E{hygienische
  Händedesinfektion} erforderlich
  (\prettyref{topic:hygienische-haendedesinfektion})!
  \item \StepHeading{Schürzen}: Einweg-Plastikschürzen bieten Schutz
  gegen Kontamination der Kleidung im Rumpf- und
  Beinbereich.
  \item \StepHeading{Schutzbrillen} sollen in allen Situationen
  getragen werden in denen mit einem Spritzen von
  schädlichen bzw. infektiösen Flüssigkeiten des Patienten
  gerechnet werden muss.
  \item \StepHeading{Atemmasken} \label{topic:atemmasken}%
  \ldots\ kommen zum Schutz vor
  luftübertragbaren Krankheiten zum Einsatz. Je nach Keim
  werden unterschiedliche Schutzmasken bzw. Schutzklassen empfohlen, welche sich
  bezüglich ihrer Schutzwirkung wesentlich unterscheiden.
  Die einfachste Form ist die altbekannte, bei Operationen verwendete,
  \T{Mund-Nasen-Maske}
  (\Abk{MNM}, \enquote*{OP-Maske}), welche für viele Situationen einen
  ausreichenden Schutz darstellt\footnote{Beim Ausbruch von
  \Abk{SARS} in Hongkong in 2003 konnten simple OP-Masken das
  Infektionsrisiko für das Personal auf 0 senken (das
  infizierte Personal trug entweder keine Maske oder einfache
  Papiermasken) \cite{Seto2003}.}.
  \T{FFP-Masken} bieten einen deutlich höhren Schutz.
  Es gibt je nach Partikelgröße, die
  ausgefiltert wird, verschiedene Klassen: FFP 1, FFP 2 und FFP
  3.
  \item \StepHeading{Schutzanzüge} kommen in speziellen Situationen zum
  Einsatz.
\end{itemize}

}{

}{}{}{}



\Layout{57}
{}
{}
{}
{%
\renewcommand{\thefootnote}{{\sffamily\alph{footnote}}}
\begin{tabularx}{\linewidth}{XXXXX}\addlinespace\toprule\addlinespace
\noindent\TabHeading{Typ}
&
\noindent\TabHeading{Mindest\-rück\-halte\-ver\-mögen des Filters}
\footnotemark[1]\footnotemark[2]
% TODO LIT: Dreller S. et.al.: RdL 66 (März 2006)
&
\noindent\TabHeading{Max. zul.
Gesamt\-leck\-age}\footnotemark[3] 
&
\noindent\TabHeading{Anwendung (Beispiele)}
\\\midrule\addlinespace
\noindent\TabSub{Mund-Nasen-Maske} & 95\PC*\footnotemark[4] & k.\,A. & Basisschutz; \Abk{MRSA}, \ldots
\\\addlinespace 
\noindent\TabSub{FFP 1} & 80\PC* & 22\PC* &
\\\addlinespace
\noindent\TabSub{FFP 2} & 94\PC* & 8\PC* & Influenza
\\\addlinespace
\noindent\TabSub{FFP 3} & 99\PC* & 2\PC* & \Abk{SARS}, offene Tuberkulose 
\\\addlinespace\bottomrule
\end{tabularx}
\footnotetext[1]{für \ce{NaCl}-Prüfaerosol}
\footnotetext[2]{Dreller S. et.al.: RdL 66 (März 2006)}
\footnotetext[3]{gem. DIN EN 149\label{fn:hyg-sm-dinen149}}
\footnotetext[4]{für Staphylococcus
aureus\label{fn:hyg-sm-stapha}}
\renewcommand{\thefootnote}{\arabic{footnote}}}
{Arten von Schutzmasken und Rückhaltevermögen}
{}


\TextSlide{Arbeitsaufteilung}
{%
Unter \prettyref{topic:einer-redet} wurde bereits erwähnt,
dass möglichst nur ein Teammitglied ein Patientengespräch
führen soll. Dies ist auch aus hygienischer Sicht ähnlich:
Derjenige, welcher direkten Patientenkontakt hat, sollte nach
Möglichkeit keinen Kontakt mit dem sauberen Material aus der
Tasche, etc. haben, um es nicht unnötig zu beschmutzen.
Benötigtes Material soll ihm von einem anderen Teammitglied
\EE{zugereicht} werden.

\Fazit{Nie mit schmutzigen Handschuhen in die Tasche greifen!
Ein Teammitglied hat Patientenkontakt, der andere reicht zu!}
}{
\begin{itemize}
  \item Material nicht unnötig kontaminieren. 
  \item Ein Teammitglied hat Patientenkontakt, Kollege reicht zu.
\end{itemize}
}{}{}{}

\Layout{20}{Umgang mit spitzen Gegenständen}
{%
\index{Kontamed}%
\index{Spitze Gegenstände!Entsorgung}%
\label{topic:nadelabwurfbehaelter}%
\noindent Spitze Gegenstände (Kanülen, Lanz\-et\-ten, \ldots )
sind sofort nach Gebrauch in den Nadelabwurfbehälter (\enquote*{Gelbe Box},
\enquote*{Konta\-med\texttrademark}, \enquote*{Sharp-Box}, \ldots )
zu ent\-sorgen!

Es gibt jedoch \EE{eine Ausnahme}: wenn aus einer Kanüle Blut
für diagnostische Zwecke entnommen werden soll, darf diese
nicht gleich verworfen werden. In der Praxis wird manchmal
der \E{Blutzuckerwert mit Blut aus der Kanüle} bestimmt, um den
Patienten einen weiteren Stich zu ersparen. Dies muss jedoch
vorher ausdrücklich kund getan werden.

\Fazit{\enquote{Wer sticht, der entsorgt!}}

\Fazit{NIE nachstopfen, NIE ausleeren. -- 
Wenn die Box voll ist, wird sie verschlossen und ausgetauscht.}
}{
\begin{itemize}
  \item Nadelabwurfbehälter.
  \begin{itemize}
    \item \enquote{Gelbe Box}
    \item Kontamed{\texttrademark}
    \item \enquote{Sharp}
  \end{itemize} 
  \item Nadel/Kanüle/Lanzette \E{sofort nach Gebrauch} entsorgen
  \item Derjenige, der sticht, entsorgt sofort!
  \item Ausnahme: Blutentnahme aus Kanüle für diagnostische
  Zwecke (Sonderfall)
\end{itemize}
}
{./images/motal-aass/kontamed_4779-00800.jpg}
{\enquote{Wer sticht, der entsorgt!}}
{Motal}

% \begin{figure}[H]
% \includegraphics[width=3.377cm,height=4.552cm]{images/AASdev20090228-img105.jpg}
% \includegraphics[width=5.71cm,height=4.289cm]{images/AASdev20090228-img106.jpg}
% \Captionof{figure}{Kontamed-Boxen, verschiedene Fabrikate.}
% \end{figure}

\TextSlide{Keimvehikel}
{%
\ldots\ wie z.\,B. Kugelschreiber, Haltegriffe, etc. sind
nicht nur für den Patienten im Sinne einer Kreuzinfektion
ein Problem, sondern auch für das Personal. Daher sind auch
in Hinblick auf den Selbstschutz die entsprechenden
Hygienemaßnahmen zu ergreifen!
}{
\begin{itemize}
  \item Kugelschreiber, Haltegriffe, Lenkrad \ldots{} 
  \item \Ca{Kreuzinfektion!}
\end{itemize}

}{}{}{}


%%%%%%%%%%%%%%%%%%%%%%%%%%%%%%%%%%%%%%%%%%%%%%%%%%%%%%%%%%%
%%%%%%%%%%%%%%%%%%%%%%%%%%%%%%%%%%%%%%%%%%%%%%%%%%%%%%%%%%%
%%%%%%%%%%%%%%%%%%%%%%%%%%%%%%%%%%%%%%%%%%%%%%%%%%%%%%%%%%%
\subsection{MRSA und andere Resistente Keime}
\index{MRSA}

Die Maßnahmen beim Transport von Patienten mit resistenten Keimen
werden unter \prettyref{MK-m:mrsa-transport} besprochen.




%%%%%%%%%%%%%%%%%%%%%%%%%%%%%%%%%%%%%%%%%%%%%%%%%%%%%%%%%%%
%%%%%%%%%%%%%%%%%%%%%%%%%%%%%%%%%%%%%%%%%%%%%%%%%%%%%%%%%%%
%%%%%%%%%%%%%%%%%%%%%%%%%%%%%%%%%%%%%%%%%%%%%%%%%%%%%%%%%%%
%%%%%%%%%%%%%%%%%%%%%%%%%%%%%%%%%%%%%%%%%%%%%%%%%%%%%%%%%%%
\subsection{Nadelstichverletzungen}
\index{Nadelstichverletzungen}
\label{topic:nadelstichverletzung}

% \MarginNo{\Z{Basierend auf: Gemeinsame Erklärung der
% Deutschen AIDS-Gesellschaft (\Abk{DAIG}) und Osterreichischen
% AIDS-Gesellschaft (\Abk{OAG}). [Post-exposure prophylaxis of HIV
% infection. German-Austrian recommendations, update
% September 2007] \emph{Dtsch Med Wochenschr}, 2009, 134 Suppl
% 1, S16-S33}\cite{OeAG2007}}

% M %%%%%%%%%%%%%%%%%%%%%%%%%%%%%%%%%%%%%%%%%%%%%%%%%%% M %
% Nadelstichverletzung
\ProcedurePrintDefault{IW49091C}

\cite{Loew200611}





\subsection{Entsorgung von Abfall}
\label{topic:entsorgung-infektioeser-abfall}
\label{topic:abfall-infektioeser-entsorgung}

Von Abfällen gehen im Wesentlichen zwei Gefahren aus,
nämlich eine \EE{Verletzungsgefahr} (Kanülen und andere
spitze Gegenstände) sowie eine \EE{Infektionsgefahr}
(z.\,B. Körperflüssigkeiten, kontaminierte Einwegartikel
und -verpackungen). 
Die \I[OeNORM S2104@\"ONORM S\,2104]{ÖNORM S\,2104} \cite{OenormS2104-2005}
unterscheidet zwischen:
\begin{enumerate}
  \item Abfälle, die \EE{weder innerhalb noch außerhalb des
  medizinischen Bereiches eine Gefahr darstellen}:
  \label{item:hyg-abfall-oenorm-gruppe1}
  
  Siedlungsabfälle und Abfälle aus Arzt-, Tierarzt- und
  Zahnarztpraxen, Hauskrankenpflege, sofern diese mit gemischten
  Siedlungsabfällen aus Haushalten und haushaltsähnlichen
  Betrieben vergleichbar sind.

  \item Abfälle, die \EE{nur innerhalb des medizinischen
  Bereiches} eine eine Infektions- oder Verletzungsgefahr
  darstellen können, jedoch nicht wie Gefährliche Abfälle
  entsorgt werden müssen:
  \label{item:hyg-abfall-oenorm-gruppe2}
  \begin{enumerate}
    \item Abfälle \EE{ohne Verletzungsgefahr} z.\,B.
    Gemische aus \EE{Wundverbänden}, Gipsverbänden,
    Stuhlwindeln, Einmalwäsche, Vorlagen, Tampons,
    Einmalartikel (z.\,B. \EE{Tupfer}, \EE{Handschuhe},
    \EE{Einmalspritzen} ohne Kanüle, \EE{Katheter},
    Infusionsgeräte ohne Dorn), restentleerte
    Urinsammelsysteme und Infusionsbeutel oder Ähnliches,
    \EE{auch wenn diese blutig sind}.
    
    Die Sammlung dieser Abfälle haben in
    ausreichend dichten Gebinden bzw. Transportbehältern zu
    erfolgen.
    \item Abfälle \EE{mit
    Verletzungsgefahr}
    z.\,B. Kanülen, Lanzetten und
    Skalpelle. Sämtliche Abfälle dieser Art müssen in
    entsprechenden Gebinden gesammelt und transportiert
    werden (\EE{Nadelabwurfbehälter},
    \FullPrettyRef{topic:nadelabwurfbehaelter}).
    \item \EE{Nassabfälle}:
    z.\,B. nicht restentleerte, mit Absaugsekreten gefüllte
    Einwegsysteme.
    
    Die Sammlung und der Transport dieser Abfälle haben in
    ausreichend dichten Gebinden bzw. Transportbehältern zu
    erfolgen.
    
    Plasma, Infusionslösungen, \EE{Blut} und \EE{Urin} sind
    unter Berücksichtigung der wasserrechtlichen
    Bestimmungen wie \EE{Abwasser} zu behandeln. Bei der
    Entleerung der Gebinde sind die entsprechenden
    Hygienemaßnahmen einzuhalten.
  \end{enumerate}
  \item Abfälle, die \EE{innerhalb und außerhalb des
  medizinischen Bereiches eine Gefahr darstellen} und in
  beiden Bereichen einer gesonderten Behandlung bedürfen
  \label{item:hyg-abfall-oenorm-gruppe3}
  
  Nicht desinfizierte mikrobiologische Kulturen und mit
  gefährlichen Erregern behafteter Abfall
  
  \Z{Diese Abfälle sind Gefahrgut im Sinne des ADR, Klasse
  6.2 (\FullPrettyRef{topic:adr}); sie sind vor der
  Abfallbereitstellung zu desinfizieren}
\end{enumerate}
Die im Rettungs- und Krankentransportdienst anfallenden
Abfällen gehören i.\,d.\,R. den Gruppen
\ref{item:hyg-abfall-oenorm-gruppe1} und
\ref{item:hyg-abfall-oenorm-gruppe2} an. Für diese ist
der Entsorgungsweg je nach Gegebenheit (Landesgesetze
/-verodnungen, Verträge mit Entsorgungsbetrieben)
unterschiedlich, die betriebsinternen
Richtlinien sind zu beachten und genau zu befolgen. Es
empfiehlt sich, kontaminiertes Material bei Übergabe des
Patienten bereits im Spital entsprechend zu entsorgen.
Ebenso ist bei umweltgefährdenden Abfällen (Batterien,
Chemikalien, \ldots) der betriebsinterne Entsorgungsweg zu
befolgen.

\nocite{KayserMikrobiologie10}
\nocite{HahnMikrobio6}
\nocite{BLIM3}
\nocite{OHCSP7}
\nocite{OeAG2007}

%%%%%%%%%%%%%%%%%%%%%%%%%%%%%%%%%%%%%%%%%%%%%%%%%%%%%%%%%%%
%%%%%%%%%%%%%%%%%%%%%%%%%%%%%%%%%%%%%%%%%%%%%%%%%%%%%%%%%%%
%%%%%%%%%%%%%%%%%%%%%%%%%%%%%%%%%%%%%%%%%%%%%%%%%%%%%%%%%%%

\begin{Literary}
  Praxisorientiertes und
  reich illustriertes Fachbuch für Sanitäter:
  {\fullcite{GruberHygieneRD1}}


  \EE{Rettungsdienstspezifisches Werk}, richtet sich
  tendenziell an Hygieneverantwortliche und -beauftrage:
  {\fullcite{WiedenmannHygieneRD1}}

 \EE{Standardwerk}:
  {\fullcite{KramerWallhaeusersPraxis1}}

 \EE{Übersichtswerk} mit
  speziellem Kapitel bezüglich Hygiene im Rettungs- und
  Krankentransport und exemplarischen
  Desinfektionsplänen:
  {\fullcite{KramerKrankenhausPraxishygiene1}}

 Werk mit Österreich-Bezug:
  {\fullcite{FlammAngewandteHygiene4}}
\end{Literary}


\ChapterBibliography

\ChapterEnd